\section{跨语言组合性与 Scaling 的错误修复}

CoT 若只是记住英语数学模板，就不应在低频语言中保持强表现。课件接着用 MGSM 展示 multilingual reasoning，再用 error analysis 与 task-spectrum diagram 讨论为什么 scaling 可能同时修复 semantic understanding、missing steps 和 arithmetic mistakes。

\subsection{Multilingual CoT：低频语言为何是强测试}

MGSM 把 250 道 GSM8K 题人工翻译为十种 typologically diverse languages，并要求模型在相应语言中读题和推理。Bengali 在预训练数据中的比例很低，因此如果模型仍能完成任务，至少说明语言理解与数学推理可以在某种程度上组合，而不只是复现大量相同语言模板。

\lecturefigure{slide-28.jpg}{Multilingual CoT：十种语言、低频 Bengali 与组合性证据}{官方课件第 28 页；Shi et al., ICLR 2023}

\begin{knowledgebox}{读图：不要只看平均柱高}
先看 benchmark 构造：相同数学问题跨语言翻译，尽量固定 reasoning difficulty；再比较 language frequency 与 accuracy 的关系。低资源语言表现好于简单频率预期，是 compositionality 的证据。它不意味着所有语言公平，也不排除翻译风格、tokenization 和共享数字符号带来的帮助。
\end{knowledgebox}

\teachervoice{讲者特别强调 Bengali 约占极少预训练数据，却出现 surprisingly good result。他把这解释为模型能组合“理解 Bengali”与“做数学”两种技能；同时在问答中承认，语言之间的推理究竟共享多少内部机制仍没有答案。}

\subsection{Scaling 修复哪些错误}

“大模型更好”过于粗糙。课件比较 PaLM 62B 与 540B 的 CoT 输出，把改进分为 semantic understanding、one-step missing 与 arithmetic 等错误类别。多类错误同时下降，说明规模并非只修复某个窄算术模块，而是在任务链的多个环节提高可靠度。

\lecturefigure{slide-29.jpg}{从 PaLM 62B 到 540B：多类 CoT 错误同时减少}{官方课件第 29 页；Wei et al., 2022}

\begin{knowledgebox}{读图：Error taxonomy 是诊断，不是完整机制}
条形图按人工分类统计较小模型的失败，以及更大模型修复了多少。若 semantic、missing step 与 arithmetic 都改善，最保守结论是 scaling 对多种 observable failure modes 有广泛作用；不能由此断言所有内部原因相同，也不能排除分类边界重叠。
\end{knowledgebox}

\subsection{任务复杂度与模型规模的二维前沿}

前面的能力前沿只画 model scale；这里再加入 task complexity。Standard prompting 覆盖容易、短路径任务，CoT 把边界推向 math word problems、symbolic reasoning 和 hard commonsense；继续增大模型又扩大可用区域。

\lecturefigure{slide-30.jpg}{Language-model scale 与 task complexity 的 CoT 能力前沿}{官方课件第 30 页；Wei et al., 2022}

\begin{knowledgebox}{读图：边界会被多个工程杠杆共同推动}
横轴是模型规模，纵轴可理解为任务所需组合步骤或难度；蓝区是 standard prompt 可做，粉区是 CoT 新增。边界并非只随参数右移：better data、fine-tuning、tools、verification 与更多 inference compute 都可能扩大区域。图是组织研究问题的概念模型，不是可直接拟合的物理相图。
\end{knowledgebox}

\begin{warningbox}{“输出更多 token”不等于“任何 filler 都有效”}
课堂问答提到 equal-length `x x x ...` 控制并不能复制 CoT 增益，先给答案再解释也通常更差。额外计算是部分直觉，但中间 token 必须携带与问题结构相关的语言信息；否则只是延长输出。
\end{warningbox}

\subsection{本章小结}

Multilingual CoT 支持技能组合而非单纯英语模板记忆，error analysis 显示 scaling 同时修复多种可观察失败，task-spectrum diagram 则把模型规模与任务复杂度放进同一框架。三者仍是行为证据，不应被夸大为内部推理机制的最终解释。

